\chapter{Review of Literature}
\label{chap:literature}

\section{Introduction}
\label{sec:lit_intro}

Visual re-identification has evolved from handcrafted feature descriptors (color histograms, local binary patterns) coupled with metric learning (Mahalanobis, KISSME), to deep convolutional and transformer representations, and recently to unsupervised and generative paradigms.

This chapter reviews foundational and modern literature relevant to this dissertation: Section~\ref{sec:lit_unsupervised_reid} explores unsupervised person and vehicle Re-ID; Section~\ref{sec:lit_part_models} analyzes part-based representations and Spatial Transformer Networks (STNs); Section~\ref{sec:lit_diffusion_models} reviews diffusion probabilistic models and score-based generative frameworks; Section~\ref{sec:lit_metric_learning} analyzes deep metric learning dynamics and proves Theorem~\ref{sec:lit_hyperspherical_saturation} regarding hyperspherical gradient saturation; and Section~\ref{sec:lit_summary} synthesizes prior literature and benchmarks state-of-the-art methods in Table~\ref{tab:lit_comparison}.

\section{Unsupervised Person and Vehicle Re-Identification}
\label{sec:lit_unsupervised_reid}

Unsupervised Re-ID approaches fall into two paradigms: \textit{unsupervised domain adaptation (UDA)} and \textit{purely unsupervised} learning.

\subsection{Unsupervised Domain Adaptation (UDA) in Re-ID}
\label{subsec:uda_reid}

Early label-free Re-ID leveraged Unsupervised Domain Adaptation (UDA)~\cite{zheng2017unlabeled,ge2020selfpaced}, adapting models from a labeled source dataset $\mathcal{D}_s = \{(\mathbf{x}_i^s, y_i^s)\}_{i=1}^{N_s}$ to an unannotated target domain $\mathcal{D}_t = \{\mathbf{x}_j^t\}_{j=1}^{N_t}$ via style-transfer GANs (e.g., CycleGAN~\cite{zheng2017unlabeled}). However, UDA relies heavily on source-target semantic alignment; severe domain gaps degrade transfer, and privacy constraints often preclude accessing labeled source data.

\subsection{Fully Unsupervised Clustering and Pseudo-Labeling Paradigms}
\label{subsec:fully_unsupervised_reid}

Modern purely unsupervised Re-ID trains exclusively on unannotated target imagery $\mathcal{X} = \{\mathbf{x}_i\}_{i=1}^N$ via iterative clustering and contrastive optimization:
\begin{enumerate}[leftmargin=*]
    \item \textbf{Self-Paced Contrastive Learning (SpCL)}~\cite{ge2020selfpaced}: Unifies source classes, target cluster centroids, and unclustered instance features in a hybrid memory bank governed by a self-paced curriculum.
    \item \textbf{Inter-Instance Contrastive Encoding (ICE)}~\cite{chen2021ice}: Exploits hard negative mining within pseudo-clusters to accelerate representation learning.
    \item \textbf{Cluster Contrast}~\cite{dai2022cluster}: Stores and updates exclusively cluster-level representations, reducing mini-batch gradient variance and improving stability.
    \item \textbf{Clustering Consensus (RLCC)}~\cite{zhang2021rlcc}: Tracks co-clustering consensus across generations to filter transient false merges and stabilize pseudo-labels.
    \item \textbf{Multi-Label and Prototypicality Formulations}: Lin~\textit{et~al.}~\cite{lin2020unsupervised} formulated multi-label classification for boundary samples, while Zhang~\textit{et~al.}~\cite{zhang2023ise} modeled emergent prototypicality.
\end{enumerate}

\subsection{Camera-Aware Proxies and Multi-Camera Domain Gaps}
\label{subsec:camera_aware_reid}

Disparate camera viewpoints, lens distortions, and illumination induce intra-identity variance exceeding inter-identity variance. To bridge this gap, Wang~\textit{et~al.}~\cite{wang2021cap} proposed \textbf{Camera-Aware Proxies (CAP)}, generating separate sub-proxies $\boldsymbol{\mu}_k^c$ per (cluster $k$, camera $c$) pair to pull features toward cross-camera intra-identity proxies. However, on large camera topologies like VeRi-776 (20 cameras)~\cite{liu2016deepveri}, unaligned DBSCAN over-segments 576 vehicles into 1,030 clusters, causing standard hard-negative mining to repel same-vehicle sub-clusters into an irreversible shattering cycle.

\subsection{Vehicle Re-Identification Benchmarks and Specific Challenges}
\label{subsec:vehicle_reid_lit}

Unlike articulated pedestrians, vehicle Re-ID presents distinct physical challenges:
\begin{itemize}[leftmargin=*]
    \item \textbf{Rigid Polyhedral Geometry}: 2D visual projections vary drastically between front, rear, and profile angles~\cite{liu2016deep}.
    \item \textbf{High Inter-Class Similarity}: Distinct vehicles frequently share identical make, model, and paint color, requiring fine-grained localization of unique inspection stickers or grilles~\cite{liu2016deepveri,lou2019veriwide}.
    \item \textbf{Viewpoint Regularization}: Spatial-temporal and orientation-invariant features~\cite{wang2017orientation,khorramshahi2019dual} disentangle viewpoint variations from identity attributes.
\end{itemize}

\section{Part-Based Representations and Spatial Transformer Networks}
\label{sec:lit_part_models}

Decomposing images into localized semantic regions provides robustness against occlusions and camera bias.

\subsection{Rigid Horizontal Stripe Partitioning}
\label{subsec:pcb_lit}

The Part-Based Convolutional Baseline (PCB)~\cite{sun2018beyond} uniformly slices feature maps into $P$ horizontal stripes (typically $P=6$). While efficient, rigid slicing fails under off-center, crouched, or occluded poses, mixing anatomical regions across adjacent stripes.

\subsection{Part-Based Pseudo-Label Refinement (PPLR)}
\label{subsec:pplr_lit}

Cho~\textit{et~al.}~\cite{cho2022pplr} proposed PPLR, using cross-agreement between global clusters and horizontal part representations to refine pseudo-labels, achieving 81.5\% mAP on Market-1501 and confirming the value of local spatial semantics.

\subsection{Spatial Transformer Networks (STNs) and Geometric Regularization}
\label{subsec:stn_lit}

Spatial Transformer Networks (STNs)~\cite{jaderberg2015spatial} provide differentiable spatial cropping via: (1) a \textbf{localization network} predicting transformation parameters $\boldsymbol{\theta} \in \mathbb{R}^{2 \times 3}$; (2) a \textbf{grid generator} computing source coordinates $(x_s, y_s)$ from regular target grid $(x_t, y_t) \in [-1, 1]^2$:
\begin{equation}
\label{eq:lit_stn_grid}
\begin{pmatrix} x_s \\ y_s \end{pmatrix} = \begin{bmatrix} \theta_{11} & \theta_{12} \\ \theta_{21} & \theta_{22} \end{bmatrix} \begin{pmatrix} x_t \\ y_t \end{pmatrix} + \begin{pmatrix} \theta_{13} \\ \theta_{23} \end{pmatrix} = \mathbf{A} \begin{pmatrix} x_t \\ y_t \end{pmatrix} + \mathbf{t}
\end{equation}
and (3) a \textbf{bilinear sampler}. Unconstrained unsupervised STNs suffer from:
\begin{itemize}[leftmargin=*]
    \item \textbf{Singular Grid Transformations}: Linear matrix $\mathbf{A}$ degenerating to $\det(\mathbf{A}) = 0$ collapses 2D receptive fields onto a 1D diagonal line (Theorem~\ref{thm:singular_grid}).
    \item \textbf{Out-of-Bounds Sampling}: Excessive translation drift pushes grids outside $[-1, 1]^2$, causing constant zero-padding.
    \item \textbf{Lightweight MLP Regularization}: In this work, canonical horizontal stripe anchors with bounded residual offsets strictly guarantee $\det(\mathbf{A}) \ge 0.196$ while reducing parameter complexity.
\end{itemize}

\section{Diffusion Probabilistic Models in Computer Vision}
\label{sec:lit_diffusion_models}

Denoising Diffusion Probabilistic Models (DDPMs)~\cite{ho2020denoising} and continuous score-based SDEs~\cite{song2020score} achieve state-of-the-art generative modeling.

\subsection{Foundations of Denoising Diffusion Probabilistic Models (DDPM)}
\label{subsec:ddpm_theory}

A standard DDPM perturbs state $\mathbf{x}_0 \sim q(\mathbf{x}_0)$ across $T$ timesteps with variance schedule $\beta_1, \dots, \beta_T$:
\begin{equation}
\label{eq:lit_ddpm_forward_step}
q(\mathbf{x}_t \mid \mathbf{x}_{t-1}) = \mathcal{N}\left(\mathbf{x}_t; \sqrt{1 - \beta_t}\mathbf{x}_{t-1}, \beta_t \mathbf{I}\right)
\end{equation}
Defining $\alpha_t = 1 - \beta_t$ and $\bar{\alpha}_t = \prod_{s=1}^t \alpha_s$, the forward marginal and denoising score matching objective satisfy:
\begin{align}
\label{eq:lit_ddpm_marginal}
q(\mathbf{x}_t \mid \mathbf{x}_0) &= \mathcal{N}\left(\mathbf{x}_t; \sqrt{\bar{\alpha}_t}\mathbf{x}_0, (1 - \bar{\alpha}_t)\mathbf{I}\right) \implies \mathbf{x}_t = \sqrt{\bar{\alpha}_t}\mathbf{x}_0 + \sqrt{1 - \bar{\alpha}_t}\boldsymbol{\epsilon} \\
\label{eq:lit_ddpm_loss}
\mathcal{L}_{\text{simple}}(\theta) &= \mathbb{E}_{t \sim \mathcal{U}\{1, T\}, \mathbf{x}_0, \boldsymbol{\epsilon} \sim \mathcal{N}(\mathbf{0}, \mathbf{I})}\left[\|\boldsymbol{\epsilon}_\theta(\mathbf{x}_t, t) - \boldsymbol{\epsilon}\|_2^2\right]
\end{align}

\subsection{Continuous Score-Based Generative Modeling (Score SDE)}
\label{subsec:score_sde}

Song~\textit{et~al.}~\cite{song2020score} unified diffusion models under continuous SDEs $\mathrm{d}\mathbf{x} = \mathbf{f}(\mathbf{x}, t)\mathrm{d}t + g(t)\mathrm{d}\mathbf{w}$ \label{eq:lit_forward_sde}, with reverse dynamics:
\begin{equation}
\label{eq:lit_reverse_sde}
\mathrm{d}\mathbf{x} = \left[\mathbf{f}(\mathbf{x}, t) - g(t)^2 \nabla_{\mathbf{x}} \log p_t(\mathbf{x})\right]\mathrm{d}t + g(t)\mathrm{d}\bar{\mathbf{w}}
\end{equation}
where Stein score estimation satisfies:
\begin{equation}
\label{eq:lit_score_relation}
\mathbf{s}_\theta(\mathbf{x}_t, t) \approx -\frac{\boldsymbol{\epsilon}_\theta(\mathbf{x}_t, t)}{\sqrt{1 - \bar{\alpha}_t}} \approx \nabla_{\mathbf{x}_t} \log p_t(\mathbf{x}_t)
\end{equation}
Crucially, \textbf{the denoiser must observe the state $\mathbf{x}_t$ being scored}. Omitting $\mathbf{x}_t$ bounds the Bayes-optimal noise loss at $1.000$ (Theorem~\ref{thm:denoiser_floor}).

\subsection{Deterministic Acceleration via Denoising Diffusion Implicit Models (DDIM)}
\label{subsec:ddim_theory}

Song~\textit{et~al.}~\cite{song2020denoising} proposed DDIM, enabling deterministic reverse transitions ($\sigma_t = 0$) over sub-sequences $\tau = [\tau_1, \dots, \tau_S]$ ($S \ll T$):
\begin{align}
\label{eq:lit_ddim_x0}
\hat{\mathbf{x}}_0 &= \frac{\mathbf{x}_{\tau_k} - \sqrt{1 - \bar{\alpha}_{\tau_k}} \, \boldsymbol{\epsilon}_\theta(\mathbf{x}_{\tau_k}, \tau_k)}{\sqrt{\bar{\alpha}_{\tau_k}}} \\
\label{eq:lit_ddim_step}
\mathbf{x}_{\tau_{k-1}} &= \sqrt{\bar{\alpha}_{\tau_{k-1}}}\hat{\mathbf{x}}_0 + \sqrt{1 - \bar{\alpha}_{\tau_{k-1}}} \, \boldsymbol{\epsilon}_\theta(\mathbf{x}_{\tau_k}, \tau_k)
\end{align}
An analytic 5-step DDIM trajectory ($\tau = [999, 749, 499, 249, 0]$) eliminates sampling latency while maintaining exact sample-to-batch index alignment.

\subsection{Diffusion Models for Structured Parameter Generation (PISL)}
\label{subsec:pisl_diffusion_lit}

While diffusion models conventionally operate on pixel space ($H \times W \times C$), Tao~\textit{et~al.}~\cite{tao2024unsupervised} applied diffusion in \textit{parameter space} over transformation parameters $\boldsymbol{\theta} \in \mathbb{R}^{P \times 2 \times 3} = \mathbb{R}^{18}$ conditioned on global feature $\mathbf{f}_g$. Our investigation identifies five optimization dimensions:
\begin{enumerate}[leftmargin=*]
    \item Omitting $\boldsymbol{\theta}_t$ from the denoiser, constraining it to unconditional prediction (Bottleneck 5, Theorem~\ref{thm:denoiser_floor}).
    \item 1000-step stochastic DDPM sampling introducing high offline clustering latency.
    \item Batch scrambling heuristics (\texttt{theta\_renewal}) disrupting image-to-parameter correspondence.
    \item Noisy checkpoint averaging (\texttt{use\_ensemble}) corrupting clean parameter estimates.
    \item Clamping vertical translation to $[0, 1]$, prohibiting upper-body crops (Bottleneck 4).
\end{enumerate}

\section{Deep Metric Learning and Hyperspherical Gradient Saturation}
\label{sec:lit_metric_learning}

Metric learning optimizes embeddings such that relative distances directly reflect identity.

\subsection{The Triplet Objective and Soft Margin Formulations}
\label{subsec:triplet_dynamics}

The foundational hard-margin triplet loss~\cite{hermans2017triplet} is $\mathcal{L}_{\text{hard\_tri}} = \max(0, \|\mathbf{f}_a - \mathbf{f}_p\|_2 - \|\mathbf{f}_a - \mathbf{f}_n\|_2 + m)$ \label{eq:lit_hard_triplet}. Hermans~\textit{et~al.}~\cite{hermans2017triplet} introduced the soft-margin formulation:
\begin{equation}
\label{eq:lit_soft_triplet}
\mathcal{L}_{\text{tri}}(\Delta) = \ln\left(1 + \exp\left(d(\mathbf{f}_a, \mathbf{f}_p) - d(\mathbf{f}_a, \mathbf{f}_n)\right)\right) = \mathrm{softplus}(\Delta), \quad \Delta = d_{ap} - d_{an}
\end{equation}

\subsection{Hyperspherical Gradient Saturation in High Dimensions}
\label{subsec:saturation_analysis}

In deep ResNet-50 backbones ($D=2048$), evaluating marginless soft-margin triplet loss on unnormalized vectors induces exponential gradient vanishing.

\begin{theorem}[Hyperspherical Gradient Saturation in High-Dimensional Spaces]
\label{sec:lit_hyperspherical_saturation}
Let $\mathbf{f}_a, \mathbf{f}_p, \mathbf{f}_n \in \mathbb{R}^D$ be unnormalized feature representations with post-BatchNorm unit variance ($\sigma_c^2 \approx 1.0$) in dimension $D = 2048$. Under the marginless soft-margin triplet objective:
\begin{equation}
\label{eq:thm_lit_triplet_loss}
\mathcal{L}_{\text{tri}}(\Delta) = \ln\left(1 + \exp(\Delta)\right), \quad \Delta = \|\mathbf{f}_a - \mathbf{f}_p\|_2 - \|\mathbf{f}_a - \mathbf{f}_n\|_2
\end{equation}
the gradient magnitude with respect to $\Delta$ satisfies:
\begin{equation}
\label{eq:thm_lit_gradient_formula}
\left| \frac{\partial \mathcal{L}_{\text{tri}}}{\partial \Delta} \right| = \frac{1}{1 + \exp(-\Delta)} = \sigma(\Delta)
\end{equation}
Because expected negative distance scales as $\mathbb{E}[\|\mathbf{f}_a - \mathbf{f}_n\|_2] \approx \sqrt{2D}\sigma_c \approx 64.0$, a separation gap of merely $\Delta \le -3.5$ attenuates the gradient to:
\begin{equation}
\label{eq:thm_lit_gradient_attenuation}
\left| \frac{\partial \mathcal{L}_{\text{tri}}}{\partial \Delta} \right| \le 0.0293
\end{equation}
reducing metric gradient updates by over $97.0\%$ and rendering unsupervised metric learning functionally dormant.
\end{theorem}

When coarse clustering contracts positive pairs to $d_{ap} \le 60.5$, $\Delta \le -3.5$ drives $\sigma(-3.5) \le 0.0293$, causing $>97\%$ gradient attenuation and freezing metric learning. For the complete derivation, see Chapter~\ref{chap:method}, Theorem~\ref{thm:triplet_saturation}.

\subsection{Resolution via Hyperspherical Projection and Temperature Scaling}
\label{subsec:resolution_hypersphere}

Optimization~P4 projects features onto $\mathbb{S}^{D-1}$ via $L_2$ normalization, strictly bounding pairwise distances:
\begin{equation}
\label{eq:lit_l2_norm}
\tilde{\mathbf{f}}_g = \frac{\mathbf{f}_g}{\|\mathbf{f}_g\|_2} \implies \tilde{d}_{ij} = \|\tilde{\mathbf{f}}_i - \tilde{\mathbf{f}}_j\|_2 = \sqrt{2 - 2\cos(\mathbf{f}_i, \mathbf{f}_j)} \in [0, 2]
\end{equation}
With margin $m = 0.3$ and inverse temperature $\beta = 2.0$:
\begin{equation}
\label{eq:lit_norm_triplet_loss}
\mathcal{L}_{\text{norm\_tri}} = \frac{1}{B}\sum_{i=1}^B \frac{1}{\beta}\ln\left(1 + \exp\left(\beta(\tilde{d}_{ap}(i) - \tilde{d}_{an}(i) + m)\right)\right)
\end{equation}
guaranteeing active gradients $|\partial \mathcal{L}_{\text{norm\_tri}} / \partial \Delta| = [1 + \exp(-\beta(\Delta + m))]^{-1} \in [0.10, 0.90]$ \label{eq:lit_active_gradient} throughout training.

\section{Literature Benchmarking and Chapter Summary}
\label{sec:lit_summary}

Table~\ref{tab:lit_comparison} presents a comprehensive comparative benchmark of representative state-of-the-art unsupervised person and vehicle re-identification methods across both Market-1501 and VeRi-776, contrasting prior literature with the published PISL results, our baseline reproduction, and our mathematically optimized framework (P1--P6).

\begin{table}[ht]
\centering
\small
\caption[Comparison of representative unsupervised Re-ID methods]{Comprehensive comparison of representative unsupervised person and vehicle Re-ID methods on Market-1501 and VeRi-776.}
\label{tab:lit_comparison}
\resizebox{\columnwidth}{!}{%
\begin{tabular}{lcccccc}
\toprule
\multirow{2}{*}{\textbf{Method}} & \multirow{2}{*}{\textbf{Venue}} & \multicolumn{2}{c}{\textbf{Market-1501}} & \multicolumn{2}{c}{\textbf{VeRi-776}} & \multirow{2}{*}{\textbf{Part / Spatial Modeling}} \\
\cmidrule(lr){3-4} \cmidrule(lr){5-6}
& & \textbf{mAP (\%)} & \textbf{Rank-1 (\%)} & \textbf{mAP (\%)} & \textbf{Rank-1 (\%)} & \\
\midrule
SpCL~\cite{ge2020selfpaced} & NeurIPS 2020 & 73.1 & 89.0 & 36.9 & 82.9 & Global Only \\
CAP~\cite{wang2021cap} & AAAI 2021 & 79.2 & 91.4 & 37.8 & 85.3 & Camera-Aware Proxies \\
RLCC~\cite{zhang2021rlcc} & CVPR 2021 & 77.7 & 90.4 & 38.3 & 84.7 & Clustering Consensus \\
ICE~\cite{chen2021ice} & ICCV 2021 & 82.3 & 93.8 & 39.1 & 86.2 & Inter-Instance Contrast \\
PPLR~\cite{cho2022pplr} & CVPR 2022 & 81.5 & 92.8 & 40.2 & 85.6 & Part Pseudo-Label Refinement \\
Cluster Contrast~\cite{dai2022cluster} & ACCV 2022 & 82.6 & 93.0 & 40.8 & 86.8 & Cluster Memory Bank \\
\midrule
PISL (Reported)~\cite{tao2024unsupervised} & IEEE TIP 2024 & 84.6 & 93.5 & 43.1 & 96.5 & STN + Parameter Diffusion \\
PISL (Reproduction Baseline) & Ours & 83.8 & 92.8 & 42.3 & 95.6 & Standard Baseline Reproduction \\
\midrule
\textbf{PISL (Optimized P1--P6)} & \textbf{Ours} & \textbf{85.6} & \textbf{94.2} & \textbf{52.4} & \textbf{96.1} & \textbf{Bounded STN + Conditioned DDIM} \\
\quad \textit{+ $k$-Reciprocal Re-Ranking}~\cite{zhong2017reranking} & \textbf{Ours} & \textbf{89.8} & \textbf{94.8} & \textbf{54.1} & \textbf{96.6} & \textbf{Manifold Post-Processing} \\
\bottomrule
\end{tabular}%
}
\end{table}

The literature demonstrates a clear evolutionary trajectory: from purely global representations (SpCL), to camera-aware proxies (CAP), to localized part refinement (PPLR), and ultimately to generative parameter-space diffusion (PISL). While the published PISL framework represented a significant conceptual leap, our mathematical investigation identifies key areas of enhancement in gradient dynamics, coordinate topology, diffusion conditioning, and metric scaling.

By deriving and applying mathematically grounded algorithmic refinements (Theorems 1 through 5, Optimizations P1 through P6), this dissertation addresses these structural bottlenecks, establishes new state-of-the-art benchmarks on both Market-1501 and VeRi-776, and provides a rigorous, publication-grade foundation for generative part discovery in unsupervised re-identification.
